% This LaTeX file was created by Metamath on 22-Jul-2024 3:28 PM.
% edited 7/17/2026
\blalldd
%\begin{restatable}{lemma}{blalldd}~\blTitle{bl.alldd}~ Universal quantification over closed time
%interval [t1 .. t2].
%\begin{align}
%\vdash ( \tau_1 \in \mathbf{T} \wedge \tau_2 \in \mathbf{T}
%  )\quad\Rightarrow\quad \vdash ( \forall \msc{x} \in \mathbf{T} ( \msc{x}
%  \in ( \tau_1 [,] \tau_2 ) \wedge \mw{\varphi} ) \leftrightarrow \forall
%  \msc{x} \in \mathbf{T} ( \tau_1 \le \msc{x} \wedge \msc{x} \le \tau_2
%  \wedge \mw{\varphi} ) )\label{eq:bl.alldd}\tag{bl.alldd}
%\end{align}
%\end{restatable}

\begin{proof}\label{proof:bl.alldd}
\begin{align}
  1 &&  & \vdash ( \tau_1 \in \mathbf{T} \wedge \tau_2 \in \mathbf{T}
      )&\text{Hyp~alldd.1}\notag%bl.alldd.1
  \\ 2 &&  & \vdash ( ( \tau_1 \in \mathbf{T} \wedge \tau_2 \in \mathbf{T} )
      \rightarrow \tau_1 \in \mathbf{T}
      )&(\mbox{\ref{eq:simpl}})\notag
  \\ 3 &&  & \vdash \tau_1 \in \mathbf{T}&(\mbox{\ref{eq:ax-mp}},1,2)\notag
  \\ 4 &&  & \vdash \mathbf{T} \subseteq
      \mathbb{R}^{\ast}&(\mbox{\ref{eq:bl.ert}})\notag
  \\ 5 &&  & \vdash ( \tau_1 \in \mathbf{T} \rightarrow \tau_1 \in
      \mathbb{R}^{\ast} )&(\mbox{\ref{eq:sseli}},4)\notag
  \\ 6 &&  & \vdash \tau_1 \in \mathbb{R}^{\ast}&(\mbox{\ref{eq:ax-mp}},3,5)\notag
  \\ 7 &&  & \vdash ( \tau_1 \in \mathbf{T} \wedge \tau_2 \in \mathbf{T}
      )&\text{Hyp~alldd.1}\notag%bl.alldd.1
  \\ 8 &&  & \vdash ( ( \tau_1 \in \mathbf{T} \wedge \tau_2 \in \mathbf{T} )
      \rightarrow \tau_2 \in \mathbf{T}
      )&(\mbox{\ref{eq:simpr}})\notag
  \\ 9 &&  & \vdash \tau_2 \in \mathbf{T}&(\mbox{\ref{eq:ax-mp}},7,8)\notag
  \\ 10 &&  & \vdash \mathbf{T} \subseteq
      \mathbb{R}^{\ast}&(\mbox{\ref{eq:bl.ert}})\notag
  \\ 11 &&  & \vdash ( \tau_2 \in \mathbf{T} \rightarrow \tau_2 \in
      \mathbb{R}^{\ast} )&(\mbox{\ref{eq:sseli}},10)\notag
  \\ 12 &&  & \vdash \tau_2 \in
      \mathbb{R}^{\ast}&(\mbox{\ref{eq:ax-mp}},9,11)\notag
  \\ 13 &&  & \vdash ( \tau_1 \in \mathbb{R}^{\ast} \wedge \tau_2 \in
      \mathbb{R}^{\ast} )&(\mbox{\ref{eq:pm3.2i}},6,12)\notag
  \\ 14 &&  & \vdash ( ( \tau_1 \in \mathbb{R}^{\ast} \wedge \tau_2 \in
      \mathbb{R}^{\ast} ) \rightarrow ( \msc{x} \in ( \tau_1
      [,] \tau_2 ) \leftrightarrow ( \msc{x} \in
      \mathbb{R}^{\ast} \wedge \tau_1 \le \msc{x} \wedge
      \msc{x} \le \tau_2 ) )
      )&(\mbox{\ref{eq:elicc1}})\notag
  \\ 15 &&  & \vdash ( \msc{x} \in ( \tau_1 [,] \tau_2 ) \leftrightarrow (
      \msc{x} \in \mathbb{R}^{\ast} \wedge \tau_1 \le \msc{x}
      \wedge \msc{x} \le \tau_2 )
      )&(\mbox{\ref{eq:ax-mp}},13,14)\notag
  \\ 16 &&  & \vdash ( ( \msc{x} \in \mathbb{R}^{\ast} \wedge \tau_1 \le \msc{x}
      \wedge \msc{x} \le \tau_2 ) \leftrightarrow ( \msc{x}
      \in \mathbb{R}^{\ast} \wedge ( \tau_1 \le \msc{x} \wedge
      \msc{x} \le \tau_2 ) )
      )&(\mbox{\ref{eq:3anass}})\notag
  \\ 17 &&  & \vdash ( \msc{x} \in ( \tau_1 [,] \tau_2 ) \leftrightarrow (
      \msc{x} \in \mathbb{R}^{\ast} \wedge ( \tau_1 \le
      \msc{x} \wedge \msc{x} \le \tau_2 ) )
      )&(\mbox{\ref{eq:bitri}},15,16)\notag
  \\ 18 &&  & \vdash ( ( \msc{x} \in ( \tau_1 [,] \tau_2 ) \wedge \mw{\varphi}
      ) \leftrightarrow ( ( \msc{x} \in \mathbb{R}^{\ast}
      \wedge ( \tau_1 \le \msc{x} \wedge \msc{x} \le \tau_2
      ) ) \wedge \mw{\varphi} )
      )&(\mbox{\ref{eq:anbi1i}},17)\notag
  \\ 19 &&  & \vdash ( ( ( \msc{x} \in \mathbb{R}^{\ast} \wedge ( \tau_1 \le
      \msc{x} \wedge \msc{x} \le \tau_2 ) ) \wedge
      \mw{\varphi} ) \leftrightarrow ( \msc{x} \in
      \mathbb{R}^{\ast} \wedge ( ( \tau_1 \le \msc{x} \wedge
      \msc{x} \le \tau_2 ) \wedge \mw{\varphi} ) )
      )&(\mbox{\ref{eq:anass}})\notag
  \\ 20 &&  & \vdash ( ( \msc{x} \in ( \tau_1 [,] \tau_2 ) \wedge \mw{\varphi}
      ) \leftrightarrow ( \msc{x} \in \mathbb{R}^{\ast} \wedge
      ( ( \tau_1 \le \msc{x} \wedge \msc{x} \le \tau_2 )
      \wedge \mw{\varphi} ) )
      )&(\mbox{\ref{eq:bitri}},18,19)\notag
  \\ 21 &&  & \vdash ( ( \tau_1 \le \msc{x} \wedge \msc{x} \le \tau_2 \wedge
      \mw{\varphi} ) \leftrightarrow ( ( \tau_1 \le \msc{x}
      \wedge \msc{x} \le \tau_2 ) \wedge \mw{\varphi} )
      )&(\mbox{\ref{eq:df-3an}})\notag
  \\ 22 &&  & \vdash ( ( \msc{x} \in \mathbb{R}^{\ast} \wedge ( \tau_1 \le \msc{x}
      \wedge \msc{x} \le \tau_2 \wedge \mw{\varphi} ) )
      \leftrightarrow ( \msc{x} \in \mathbb{R}^{\ast} \wedge (
      ( \tau_1 \le \msc{x} \wedge \msc{x} \le \tau_2 )
      \wedge \mw{\varphi} ) )
      )&(\mbox{\ref{eq:anbi2i}},21)\notag
  \\ 23 &&  & \vdash ( ( \msc{x} \in ( \tau_1 [,] \tau_2 ) \wedge \mw{\varphi}
      ) \leftrightarrow ( \msc{x} \in \mathbb{R}^{\ast} \wedge
      ( \tau_1 \le \msc{x} \wedge \msc{x} \le \tau_2 \wedge
      \mw{\varphi} ) )
      )&(\mbox{\ref{eq:bitr4i}},20,22)\notag
  \\ 24 &&  & \vdash ( \forall \msc{x} \in \mathbf{T} ( \msc{x} \in ( \tau_1
      [,] \tau_2 ) \wedge \mw{\varphi} ) \leftrightarrow
      \forall \msc{x} \in \mathbf{T} ( \msc{x} \in
      \mathbb{R}^{\ast} \wedge ( \tau_1 \le \msc{x} \wedge
      \msc{x} \le \tau_2 \wedge \mw{\varphi} ) )
      )&(\mbox{\ref{eq:ralbii}},23)\notag
  \\ 25 &&  & \vdash \mathbf{T} \subseteq
      \mathbb{R}^{\ast}&(\mbox{\ref{eq:bl.ert}})\notag
  \\ 26 &&  & \vdash ( \forall \msc{x} \in \mathbf{T} ( \tau_1 \le \msc{x}
      \wedge \msc{x} \le \tau_2 \wedge \mw{\varphi} )
      \leftrightarrow \forall \msc{x} \in \mathbf{T} (
      \msc{x} \in \mathbb{R}^{\ast} \wedge ( \tau_1 \le
      \msc{x} \wedge \msc{x} \le \tau_2 \wedge \mw{\varphi}
      ) ) )&(\mbox{\ref{eq:bl.ralssi}},25)\notag
  \\ 27 &&  & \vdash ( \forall \msc{x} \in \mathbf{T} ( \msc{x} \in ( \tau_1
      [,] \tau_2 ) \wedge \mw{\varphi} ) \leftrightarrow
      \forall \msc{x} \in \mathbf{T} ( \tau_1 \le \msc{x}
      \wedge \msc{x} \le \tau_2 \wedge \mw{\varphi} )
      )&(\mbox{\ref{eq:bitr4i}},24,26)\notag
\end{align}
\end{proof}

\blallcd
%\begin{restatable}{lemma}{blallcd}~\blTitle{bl.allcd}~ Universal quantification over open-left
%time interval [t1 ,. t2].
%\begin{align}
%\vdash ( \tau_1 \in \mathbf{T} \wedge \tau_2 \in \mathbf{T}
%    )\quad\Rightarrow\quad \vdash ( \forall \msc{x} \in \mathbf{T} ( \msc{x}
%    \in ( \tau_1 (,] \tau_2 ) \wedge \mw{\varphi} ) \leftrightarrow \forall
%    \msc{x} \in \mathbf{T} ( \tau_1 < \msc{x} \wedge \msc{x} \le \tau_2 \wedge
%    \mw{\varphi} ) )\label{eq:bl.allcd}\tag{bl.allcd}
%\end{align}
%\end{restatable}

\begin{proof}\label{proof:bl.allcd}
\begin{align}
  1 &&  & \vdash ( \tau_1 \in \mathbf{T} \wedge \tau_2 \in \mathbf{T}
      )&\text{Hyp~allcd.1}\notag%bl.allcd.1
  \\ 2 &&  & \vdash ( ( \tau_1 \in \mathbf{T} \wedge \tau_2 \in \mathbf{T} )
      \rightarrow \tau_1 \in \mathbf{T}
      )&(\mbox{\ref{eq:simpl}})\notag
  \\ 3 &&  & \vdash \tau_1 \in \mathbf{T}&(\mbox{\ref{eq:ax-mp}},1,2)\notag
  \\ 4 &&  & \vdash \mathbf{T} \subseteq
      \mathbb{R}^{\ast}&(\mbox{\ref{eq:bl.ert}})\notag
  \\ 5 &&  & \vdash ( \tau_1 \in \mathbf{T} \rightarrow \tau_1 \in
      \mathbb{R}^{\ast} )&(\mbox{\ref{eq:sseli}},4)\notag
  \\ 6 &&  & \vdash \tau_1 \in \mathbb{R}^{\ast}&(\mbox{\ref{eq:ax-mp}},3,5)\notag
  \\ 7 &&  & \vdash ( \tau_1 \in \mathbf{T} \wedge \tau_2 \in \mathbf{T}
      )&\text{Hyp~allcd.1}\notag%bl.allcd.1
  \\ 8 &&  & \vdash ( ( \tau_1 \in \mathbf{T} \wedge \tau_2 \in \mathbf{T} )
      \rightarrow \tau_2 \in \mathbf{T}
      )&(\mbox{\ref{eq:simpr}})\notag
  \\ 9 &&  & \vdash \tau_2 \in \mathbf{T}&(\mbox{\ref{eq:ax-mp}},7,8)\notag
  \\ 10 &&  & \vdash \mathbf{T} \subseteq
      \mathbb{R}^{\ast}&(\mbox{\ref{eq:bl.ert}})\notag
  \\ 11 &&  & \vdash ( \tau_2 \in \mathbf{T} \rightarrow \tau_2 \in
      \mathbb{R}^{\ast} )&(\mbox{\ref{eq:sseli}},10)\notag
  \\ 12 &&  & \vdash \tau_2 \in
      \mathbb{R}^{\ast}&(\mbox{\ref{eq:ax-mp}},9,11)\notag
  \\ 13 &&  & \vdash ( \tau_1 \in \mathbb{R}^{\ast} \wedge \tau_2 \in
      \mathbb{R}^{\ast} )&(\mbox{\ref{eq:pm3.2i}},6,12)\notag
  \\ 14 &&  & \vdash ( ( \tau_1 \in \mathbb{R}^{\ast} \wedge \tau_2 \in
      \mathbb{R}^{\ast} ) \rightarrow ( \msc{x} \in ( \tau_1
      (,] \tau_2 ) \leftrightarrow ( \msc{x} \in
      \mathbb{R}^{\ast} \wedge \tau_1 < \msc{x} \wedge \msc{x}
      \le \tau_2 ) ) )&(\mbox{\ref{eq:elioc1}})\notag
  \\ 15 &&  & \vdash ( \msc{x} \in ( \tau_1 (,] \tau_2 ) \leftrightarrow (
      \msc{x} \in \mathbb{R}^{\ast} \wedge \tau_1 < \msc{x}
      \wedge \msc{x} \le \tau_2 )
      )&(\mbox{\ref{eq:ax-mp}},13,14)\notag
  \\ 16 &&  & \vdash ( ( \msc{x} \in \mathbb{R}^{\ast} \wedge \tau_1 < \msc{x}
      \wedge \msc{x} \le \tau_2 ) \leftrightarrow ( \msc{x}
      \in \mathbb{R}^{\ast} \wedge ( \tau_1 < \msc{x} \wedge
      \msc{x} \le \tau_2 ) )
      )&(\mbox{\ref{eq:3anass}})\notag
  \\ 17 &&  & \vdash ( \msc{x} \in ( \tau_1 (,] \tau_2 ) \leftrightarrow (
      \msc{x} \in \mathbb{R}^{\ast} \wedge ( \tau_1 < \msc{x}
      \wedge \msc{x} \le \tau_2 ) )
      )&(\mbox{\ref{eq:bitri}},15,16)\notag
  \\ 18 &&  & \vdash ( ( \msc{x} \in ( \tau_1 (,] \tau_2 ) \wedge \mw{\varphi}
      ) \leftrightarrow ( ( \msc{x} \in \mathbb{R}^{\ast}
      \wedge ( \tau_1 < \msc{x} \wedge \msc{x} \le \tau_2 )
      ) \wedge \mw{\varphi} )
      )&(\mbox{\ref{eq:anbi1i}},17)\notag
  \\ 19 &&  & \vdash ( ( ( \msc{x} \in \mathbb{R}^{\ast} \wedge ( \tau_1 < \msc{x}
      \wedge \msc{x} \le \tau_2 ) ) \wedge \mw{\varphi} )
      \leftrightarrow ( \msc{x} \in \mathbb{R}^{\ast} \wedge (
      ( \tau_1 < \msc{x} \wedge \msc{x} \le \tau_2 ) \wedge
      \mw{\varphi} ) ) )&(\mbox{\ref{eq:anass}})\notag
  \\ 20 &&  & \vdash ( ( \msc{x} \in ( \tau_1 (,] \tau_2 ) \wedge \mw{\varphi}
      ) \leftrightarrow ( \msc{x} \in \mathbb{R}^{\ast} \wedge
      ( ( \tau_1 < \msc{x} \wedge \msc{x} \le \tau_2 )
      \wedge \mw{\varphi} ) )
      )&(\mbox{\ref{eq:bitri}},18,19)\notag
  \\ 21 &&  & \vdash ( ( \tau_1 < \msc{x} \wedge \msc{x} \le \tau_2 \wedge
      \mw{\varphi} ) \leftrightarrow ( ( \tau_1 < \msc{x}
      \wedge \msc{x} \le \tau_2 ) \wedge \mw{\varphi} )
      )&(\mbox{\ref{eq:df-3an}})\notag
  \\ 22 &&  & \vdash ( ( \msc{x} \in \mathbb{R}^{\ast} \wedge ( \tau_1 < \msc{x}
      \wedge \msc{x} \le \tau_2 \wedge \mw{\varphi} ) )
      \leftrightarrow ( \msc{x} \in \mathbb{R}^{\ast} \wedge (
      ( \tau_1 < \msc{x} \wedge \msc{x} \le \tau_2 ) \wedge
      \mw{\varphi} ) ) )&(\mbox{\ref{eq:anbi2i}},21)\notag
  \\ 23 &&  & \vdash ( ( \msc{x} \in ( \tau_1 (,] \tau_2 ) \wedge \mw{\varphi}
      ) \leftrightarrow ( \msc{x} \in \mathbb{R}^{\ast} \wedge
      ( \tau_1 < \msc{x} \wedge \msc{x} \le \tau_2 \wedge
      \mw{\varphi} ) )
      )&(\mbox{\ref{eq:bitr4i}},20,22)\notag
  \\ 24 &&  & \vdash ( \forall \msc{x} \in \mathbf{T} ( \msc{x} \in ( \tau_1
      (,] \tau_2 ) \wedge \mw{\varphi} ) \leftrightarrow
      \forall \msc{x} \in \mathbf{T} ( \msc{x} \in
      \mathbb{R}^{\ast} \wedge ( \tau_1 < \msc{x} \wedge
      \msc{x} \le \tau_2 \wedge \mw{\varphi} ) )
      )&(\mbox{\ref{eq:ralbii}},23)\notag
  \\ 25 &&  & \vdash \mathbf{T} \subseteq
      \mathbb{R}^{\ast}&(\mbox{\ref{eq:bl.ert}})\notag
  \\ 26 &&  & \vdash ( \forall \msc{x} \in \mathbf{T} ( \tau_1 < \msc{x} \wedge
      \msc{x} \le \tau_2 \wedge \mw{\varphi} )
      \leftrightarrow \forall \msc{x} \in \mathbf{T} (
      \msc{x} \in \mathbb{R}^{\ast} \wedge ( \tau_1 < \msc{x}
      \wedge \msc{x} \le \tau_2 \wedge \mw{\varphi} ) )
      )&(\mbox{\ref{eq:bl.ralssi}},25)\notag
  \\ 27 &&  & \vdash ( \forall \msc{x} \in \mathbf{T} ( \msc{x} \in ( \tau_1
      (,] \tau_2 ) \wedge \mw{\varphi} ) \leftrightarrow
      \forall \msc{x} \in \mathbf{T} ( \tau_1 < \msc{x}
      \wedge \msc{x} \le \tau_2 \wedge \mw{\varphi} )
      )&(\mbox{\ref{eq:bitr4i}},24,26)\notag
\end{align}
\end{proof}

\blalldc
%\begin{restatable}{lemma}{blalldc}~\blTitle{bl.alldc}~ Universal quantification over open-right
%time interval [t1 ., t2].
%\begin{align}
%\vdash ( \tau_1 \in \mathbf{T} \wedge \tau_2 \in \mathbf{T}
%    )\quad\Rightarrow\quad \vdash ( \forall \msc{x} \in \mathbf{T} ( \msc{x}
%    \in ( \tau_1 [,) \tau_2 ) \wedge \mw{\varphi} ) \leftrightarrow \forall
%    \msc{x} \in \mathbf{T} ( \tau_1 \le \msc{x} \wedge \msc{x} < \tau_2 \wedge
%    \mw{\varphi} ) )\label{eq:bl.alldc}\tag{bl.alldc}
%\end{align}
%\end{restatable}

\begin{proof}\label{proof:bl.alldc}
\begin{align}
  1 &&  & \vdash ( \tau_1 \in \mathbf{T} \wedge \tau_2 \in \mathbf{T}
      )&\text{Hyp~alldc.1}\notag%bl.alldc.1
  \\ 2 &&  & \vdash ( ( \tau_1 \in \mathbf{T} \wedge \tau_2 \in \mathbf{T} )
      \rightarrow \tau_1 \in \mathbf{T}
      )&(\mbox{\ref{eq:simpl}})\notag
  \\ 3 &&  & \vdash \tau_1 \in \mathbf{T}&(\mbox{\ref{eq:ax-mp}},1,2)\notag
  \\ 4 &&  & \vdash \mathbf{T} \subseteq
      \mathbb{R}^{\ast}&(\mbox{\ref{eq:bl.ert}})\notag
  \\ 5 &&  & \vdash ( \tau_1 \in \mathbf{T} \rightarrow \tau_1 \in
      \mathbb{R}^{\ast} )&(\mbox{\ref{eq:sseli}},4)\notag
  \\ 6 &&  & \vdash \tau_1 \in \mathbb{R}^{\ast}&(\mbox{\ref{eq:ax-mp}},3,5)\notag
  \\ 7 &&  & \vdash ( \tau_1 \in \mathbf{T} \wedge \tau_2 \in \mathbf{T}
      )&\text{Hyp~alldc.1}\notag%bl.alldc.1
  \\ 8 &&  & \vdash ( ( \tau_1 \in \mathbf{T} \wedge \tau_2 \in \mathbf{T} )
      \rightarrow \tau_2 \in \mathbf{T}
      )&(\mbox{\ref{eq:simpr}})\notag
  \\ 9 &&  & \vdash \tau_2 \in \mathbf{T}&(\mbox{\ref{eq:ax-mp}},7,8)\notag
  \\ 10 &&  & \vdash \mathbf{T} \subseteq
      \mathbb{R}^{\ast}&(\mbox{\ref{eq:bl.ert}})\notag
  \\ 11 &&  & \vdash ( \tau_2 \in \mathbf{T} \rightarrow \tau_2 \in
      \mathbb{R}^{\ast} )&(\mbox{\ref{eq:sseli}},10)\notag
  \\ 12 &&  & \vdash \tau_2 \in
      \mathbb{R}^{\ast}&(\mbox{\ref{eq:ax-mp}},9,11)\notag
  \\ 13 &&  & \vdash ( \tau_1 \in \mathbb{R}^{\ast} \wedge \tau_2 \in
      \mathbb{R}^{\ast} )&(\mbox{\ref{eq:pm3.2i}},6,12)\notag
  \\ 14 &&  & \vdash ( ( \tau_1 \in \mathbb{R}^{\ast} \wedge \tau_2 \in
      \mathbb{R}^{\ast} ) \rightarrow ( \msc{x} \in ( \tau_1
      [,) \tau_2 ) \leftrightarrow ( \msc{x} \in
      \mathbb{R}^{\ast} \wedge \tau_1 \le \msc{x} \wedge
      \msc{x} < \tau_2 ) ) )&(\mbox{\ref{eq:elico1}})\notag
  \\ 15 &&  & \vdash ( \msc{x} \in ( \tau_1 [,) \tau_2 ) \leftrightarrow (
      \msc{x} \in \mathbb{R}^{\ast} \wedge \tau_1 \le \msc{x}
      \wedge \msc{x} < \tau_2 )
      )&(\mbox{\ref{eq:ax-mp}},13,14)\notag
  \\ 16 &&  & \vdash ( ( \msc{x} \in \mathbb{R}^{\ast} \wedge \tau_1 \le \msc{x}
      \wedge \msc{x} < \tau_2 ) \leftrightarrow ( \msc{x}
      \in \mathbb{R}^{\ast} \wedge ( \tau_1 \le \msc{x} \wedge
      \msc{x} < \tau_2 ) ) )&(\mbox{\ref{eq:3anass}})\notag
  \\ 17 &&  & \vdash ( \msc{x} \in ( \tau_1 [,) \tau_2 ) \leftrightarrow (
      \msc{x} \in \mathbb{R}^{\ast} \wedge ( \tau_1 \le
      \msc{x} \wedge \msc{x} < \tau_2 ) )
      )&(\mbox{\ref{eq:bitri}},15,16)\notag
  \\ 18 &&  & \vdash ( ( \msc{x} \in ( \tau_1 [,) \tau_2 ) \wedge \mw{\varphi}
      ) \leftrightarrow ( ( \msc{x} \in \mathbb{R}^{\ast}
      \wedge ( \tau_1 \le \msc{x} \wedge \msc{x} < \tau_2 )
      ) \wedge \mw{\varphi} )
      )&(\mbox{\ref{eq:anbi1i}},17)\notag
  \\ 19 &&  & \vdash ( ( ( \msc{x} \in \mathbb{R}^{\ast} \wedge ( \tau_1 \le
      \msc{x} \wedge \msc{x} < \tau_2 ) ) \wedge
      \mw{\varphi} ) \leftrightarrow ( \msc{x} \in
      \mathbb{R}^{\ast} \wedge ( ( \tau_1 \le \msc{x} \wedge
      \msc{x} < \tau_2 ) \wedge \mw{\varphi} ) )
      )&(\mbox{\ref{eq:anass}})\notag
  \\ 20 &&  & \vdash ( ( \msc{x} \in ( \tau_1 [,) \tau_2 ) \wedge \mw{\varphi}
      ) \leftrightarrow ( \msc{x} \in \mathbb{R}^{\ast} \wedge
      ( ( \tau_1 \le \msc{x} \wedge \msc{x} < \tau_2 )
      \wedge \mw{\varphi} ) )
      )&(\mbox{\ref{eq:bitri}},18,19)\notag
  \\ 21 &&  & \vdash ( ( \tau_1 \le \msc{x} \wedge \msc{x} < \tau_2 \wedge
      \mw{\varphi} ) \leftrightarrow ( ( \tau_1 \le \msc{x}
      \wedge \msc{x} < \tau_2 ) \wedge \mw{\varphi} )
      )&(\mbox{\ref{eq:df-3an}})\notag
  \\ 22 &&  & \vdash ( ( \msc{x} \in \mathbb{R}^{\ast} \wedge ( \tau_1 \le \msc{x}
      \wedge \msc{x} < \tau_2 \wedge \mw{\varphi} ) )
      \leftrightarrow ( \msc{x} \in \mathbb{R}^{\ast} \wedge (
      ( \tau_1 \le \msc{x} \wedge \msc{x} < \tau_2 ) \wedge
      \mw{\varphi} ) ) )&(\mbox{\ref{eq:anbi2i}},21)\notag
  \\ 23 &&  & \vdash ( ( \msc{x} \in ( \tau_1 [,) \tau_2 ) \wedge \mw{\varphi}
      ) \leftrightarrow ( \msc{x} \in \mathbb{R}^{\ast} \wedge
      ( \tau_1 \le \msc{x} \wedge \msc{x} < \tau_2 \wedge
      \mw{\varphi} ) )
      )&(\mbox{\ref{eq:bitr4i}},20,22)\notag
  \\ 24 &&  & \vdash ( \forall \msc{x} \in \mathbf{T} ( \msc{x} \in ( \tau_1
      [,) \tau_2 ) \wedge \mw{\varphi} ) \leftrightarrow
      \forall \msc{x} \in \mathbf{T} ( \msc{x} \in
      \mathbb{R}^{\ast} \wedge ( \tau_1 \le \msc{x} \wedge
      \msc{x} < \tau_2 \wedge \mw{\varphi} ) )
      )&(\mbox{\ref{eq:ralbii}},23)\notag
  \\ 25 &&  & \vdash \mathbf{T} \subseteq
      \mathbb{R}^{\ast}&(\mbox{\ref{eq:bl.ert}})\notag
  \\ 26 &&  & \vdash ( \forall \msc{x} \in \mathbf{T} ( \tau_1 \le \msc{x}
      \wedge \msc{x} < \tau_2 \wedge \mw{\varphi} )
      \leftrightarrow \forall \msc{x} \in \mathbf{T} (
      \msc{x} \in \mathbb{R}^{\ast} \wedge ( \tau_1 \le
      \msc{x} \wedge \msc{x} < \tau_2 \wedge \mw{\varphi} )
      ) )&(\mbox{\ref{eq:bl.ralssi}},25)\notag
  \\ 27 &&  & \vdash ( \forall \msc{x} \in \mathbf{T} ( \msc{x} \in ( \tau_1
      [,) \tau_2 ) \wedge \mw{\varphi} ) \leftrightarrow
      \forall \msc{x} \in \mathbf{T} ( \tau_1 \le \msc{x}
      \wedge \msc{x} < \tau_2 \wedge \mw{\varphi} )
      )&(\mbox{\ref{eq:bitr4i}},24,26)\notag
\end{align}
\end{proof}

\blallcc
%\begin{restatable}{lemma}{blallcc}~\blTitle{bl.allcc}~ Universal quantification over open time
%interval [t1 ,, t2].
%\begin{align}
%\vdash ( \tau_1 \in \mathbf{T} \wedge \tau_2 \in \mathbf{T}
%    )\quad\Rightarrow\quad \vdash ( \forall \msc{x} \in \mathbf{T} ( \msc{x}
%    \in ( \tau_1 (,) \tau_2 ) \wedge \mw{\varphi} ) \leftrightarrow \forall
%    \msc{x} \in \mathbf{T} ( \tau_1 < \msc{x} \wedge \msc{x} < \tau_2 \wedge
%    \mw{\varphi} ) )\label{eq:bl.allcc}\tag{bl.allcc}
%\end{align}
%\end{restatable}

\begin{proof}\label{proof:bl.allcc}
\begin{align}
  1 &&  & \vdash ( \tau_1 \in \mathbf{T} \wedge \tau_2 \in \mathbf{T}
      )&\text{Hyp~allcc.1}\notag%bl.allcc.1
  \\ 2 &&  & \vdash ( ( \tau_1 \in \mathbf{T} \wedge \tau_2 \in \mathbf{T} )
      \rightarrow \tau_1 \in \mathbf{T}
      )&(\mbox{\ref{eq:simpl}})\notag
  \\ 3 &&  & \vdash \tau_1 \in \mathbf{T}&(\mbox{\ref{eq:ax-mp}},1,2)\notag
  \\ 4 &&  & \vdash \mathbf{T} \subseteq
      \mathbb{R}^{\ast}&(\mbox{\ref{eq:bl.ert}})\notag
  \\ 5 &&  & \vdash ( \tau_1 \in \mathbf{T} \rightarrow \tau_1 \in
      \mathbb{R}^{\ast} )&(\mbox{\ref{eq:sseli}},4)\notag
  \\ 6 &&  & \vdash \tau_1 \in \mathbb{R}^{\ast}&(\mbox{\ref{eq:ax-mp}},3,5)\notag
  \\ 7 &&  & \vdash ( \tau_1 \in \mathbf{T} \wedge \tau_2 \in \mathbf{T}
      )&\text{Hyp~allcc.1}\notag%bl.allcc.1
  \\ 8 &&  & \vdash ( ( \tau_1 \in \mathbf{T} \wedge \tau_2 \in \mathbf{T} )
      \rightarrow \tau_2 \in \mathbf{T}
      )&(\mbox{\ref{eq:simpr}})\notag
  \\ 9 &&  & \vdash \tau_2 \in \mathbf{T}&(\mbox{\ref{eq:ax-mp}},7,8)\notag
  \\ 10 &&  & \vdash \mathbf{T} \subseteq
      \mathbb{R}^{\ast}&(\mbox{\ref{eq:bl.ert}})\notag
  \\ 11 &&  & \vdash ( \tau_2 \in \mathbf{T} \rightarrow \tau_2 \in
      \mathbb{R}^{\ast} )&(\mbox{\ref{eq:sseli}},10)\notag
  \\ 12 &&  & \vdash \tau_2 \in
      \mathbb{R}^{\ast}&(\mbox{\ref{eq:ax-mp}},9,11)\notag
  \\ 13 &&  & \vdash ( \tau_1 \in \mathbb{R}^{\ast} \wedge \tau_2 \in
      \mathbb{R}^{\ast} )&(\mbox{\ref{eq:pm3.2i}},6,12)\notag
  \\ 14 &&  & \vdash ( ( \tau_1 \in \mathbb{R}^{\ast} \wedge \tau_2 \in
      \mathbb{R}^{\ast} ) \rightarrow ( \msc{x} \in ( \tau_1
      (,) \tau_2 ) \leftrightarrow ( \msc{x} \in
      \mathbb{R}^{\ast} \wedge \tau_1 < \msc{x} \wedge \msc{x}
      < \tau_2 ) ) )&(\mbox{\ref{eq:elioo1}})\notag
  \\ 15 &&  & \vdash ( \msc{x} \in ( \tau_1 (,) \tau_2 ) \leftrightarrow (
      \msc{x} \in \mathbb{R}^{\ast} \wedge \tau_1 < \msc{x}
      \wedge \msc{x} < \tau_2 )
      )&(\mbox{\ref{eq:ax-mp}},13,14)\notag
  \\ 16 &&  & \vdash ( ( \msc{x} \in \mathbb{R}^{\ast} \wedge \tau_1 < \msc{x}
      \wedge \msc{x} < \tau_2 ) \leftrightarrow ( \msc{x}
      \in \mathbb{R}^{\ast} \wedge ( \tau_1 < \msc{x} \wedge
      \msc{x} < \tau_2 ) ) )&(\mbox{\ref{eq:3anass}})\notag
  \\ 17 &&  & \vdash ( \msc{x} \in ( \tau_1 (,) \tau_2 ) \leftrightarrow (
      \msc{x} \in \mathbb{R}^{\ast} \wedge ( \tau_1 < \msc{x}
      \wedge \msc{x} < \tau_2 ) )
      )&(\mbox{\ref{eq:bitri}},15,16)\notag
  \\ 18 &&  & \vdash ( ( \msc{x} \in ( \tau_1 (,) \tau_2 ) \wedge \mw{\varphi}
      ) \leftrightarrow ( ( \msc{x} \in \mathbb{R}^{\ast}
      \wedge ( \tau_1 < \msc{x} \wedge \msc{x} < \tau_2 ) )
      \wedge \mw{\varphi} )
      )&(\mbox{\ref{eq:anbi1i}},17)\notag
  \\ 19 &&  & \vdash ( ( ( \msc{x} \in \mathbb{R}^{\ast} \wedge ( \tau_1 < \msc{x}
      \wedge \msc{x} < \tau_2 ) ) \wedge \mw{\varphi} )
      \leftrightarrow ( \msc{x} \in \mathbb{R}^{\ast} \wedge (
      ( \tau_1 < \msc{x} \wedge \msc{x} < \tau_2 ) \wedge
      \mw{\varphi} ) ) )&(\mbox{\ref{eq:anass}})\notag
  \\ 20 &&  & \vdash ( ( \msc{x} \in ( \tau_1 (,) \tau_2 ) \wedge \mw{\varphi}
      ) \leftrightarrow ( \msc{x} \in \mathbb{R}^{\ast} \wedge
      ( ( \tau_1 < \msc{x} \wedge \msc{x} < \tau_2 ) \wedge
      \mw{\varphi} ) )
      )&(\mbox{\ref{eq:bitri}},18,19)\notag
  \\ 21 &&  & \vdash ( ( \tau_1 < \msc{x} \wedge \msc{x} < \tau_2 \wedge
      \mw{\varphi} ) \leftrightarrow ( ( \tau_1 < \msc{x}
      \wedge \msc{x} < \tau_2 ) \wedge \mw{\varphi} )
      )&(\mbox{\ref{eq:df-3an}})\notag
  \\ 22 &&  & \vdash ( ( \msc{x} \in \mathbb{R}^{\ast} \wedge ( \tau_1 < \msc{x}
      \wedge \msc{x} < \tau_2 \wedge \mw{\varphi} ) )
      \leftrightarrow ( \msc{x} \in \mathbb{R}^{\ast} \wedge (
      ( \tau_1 < \msc{x} \wedge \msc{x} < \tau_2 ) \wedge
      \mw{\varphi} ) ) )&(\mbox{\ref{eq:anbi2i}},21)\notag
  \\ 23 &&  & \vdash ( ( \msc{x} \in ( \tau_1 (,) \tau_2 ) \wedge \mw{\varphi}
      ) \leftrightarrow ( \msc{x} \in \mathbb{R}^{\ast} \wedge
      ( \tau_1 < \msc{x} \wedge \msc{x} < \tau_2 \wedge
      \mw{\varphi} ) )
      )&(\mbox{\ref{eq:bitr4i}},20,22)\notag
  \\ 24 &&  & \vdash ( \forall \msc{x} \in \mathbf{T} ( \msc{x} \in ( \tau_1
      (,) \tau_2 ) \wedge \mw{\varphi} ) \leftrightarrow
      \forall \msc{x} \in \mathbf{T} ( \msc{x} \in
      \mathbb{R}^{\ast} \wedge ( \tau_1 < \msc{x} \wedge
      \msc{x} < \tau_2 \wedge \mw{\varphi} ) )
      )&(\mbox{\ref{eq:ralbii}},23)\notag
  \\ 25 &&  & \vdash \mathbf{T} \subseteq
      \mathbb{R}^{\ast}&(\mbox{\ref{eq:bl.ert}})\notag
  \\ 26 &&  & \vdash ( \forall \msc{x} \in \mathbf{T} ( \tau_1 < \msc{x} \wedge
      \msc{x} < \tau_2 \wedge \mw{\varphi} )
      \leftrightarrow \forall \msc{x} \in \mathbf{T} (
      \msc{x} \in \mathbb{R}^{\ast} \wedge ( \tau_1 < \msc{x}
      \wedge \msc{x} < \tau_2 \wedge \mw{\varphi} ) )
      )&(\mbox{\ref{eq:bl.ralssi}},25)\notag
  \\ 27 &&  & \vdash ( \forall \msc{x} \in \mathbf{T} ( \msc{x} \in ( \tau_1
      (,) \tau_2 ) \wedge \mw{\varphi} ) \leftrightarrow
      \forall \msc{x} \in \mathbf{T} ( \tau_1 < \msc{x}
      \wedge \msc{x} < \tau_2 \wedge \mw{\varphi} )
      )&(\mbox{\ref{eq:bitr4i}},24,26)\notag
\end{align}
\end{proof}

\blexdd
%\begin{restatable}{lemma}{blexdd}~\blTitle{bl.exdd}~ Existential quantification over closed
%interval [t1 .. t2].
%\begin{align}
%\vdash ( \exists \msc{x} ( \msc{x} \in ( \tau_1 [,] \tau_2 ) \wedge
%    \mw{\varphi} ) \leftrightarrow \lnot \forall \msc{x} \lnot ( \msc{x} \in (
%    \tau_1 [,] \tau_2 ) \wedge \mw{\varphi} ) )\label{eq:bl.exdd}\tag{bl.exdd}
%\end{align}
%\end{restatable}

\begin{proof}\label{proof:bl.exdd}
\begin{align}
  1 &&  & \vdash ( \exists \msc{x} ( \msc{x} \in ( \tau_1 [,] \tau_2 ) \wedge
      \mw{\varphi} ) \leftrightarrow \lnot \forall \msc{x}
      \lnot ( \msc{x} \in ( \tau_1 [,] \tau_2 ) \wedge
      \mw{\varphi} ) )&(\mbox{\ref{eq:df-ex}})\notag
\end{align}
\end{proof}

\blexcc
%\begin{restatable}{lemma}{blexcc}~\blTitle{bl.excc}~ Existential quantification over open
%interval [t1 ,, t2].
%\begin{align}
%\vdash ( \exists \msc{x} ( \msc{x} \in ( \tau_1 (,) \tau_2 ) \wedge
%    \mw{\varphi} ) \leftrightarrow \lnot \forall \msc{x} \lnot ( \msc{x} \in (
%    \tau_1 (,) \tau_2 ) \wedge \mw{\varphi} ) )\label{eq:bl.excc}\tag{bl.excc}
%\end{align}
%\end{restatable}

\begin{proof}\label{proof:bl.excc}
\begin{align}
  1 &&  & \vdash ( \exists \msc{x} ( \msc{x} \in ( \tau_1 (,) \tau_2 ) \wedge
      \mw{\varphi} ) \leftrightarrow \lnot \forall \msc{x}
      \lnot ( \msc{x} \in ( \tau_1 (,) \tau_2 ) \wedge
      \mw{\varphi} ) )&(\mbox{\ref{eq:df-ex}})\notag
\end{align}
\end{proof}

\blexcd
%\begin{restatable}{lemma}{blexcd}~\blTitle{bl.excd}~ Existential quantification over open-left
%interval [t1 ,. t2].
%\begin{align}
%\vdash ( \exists \msc{x} ( \msc{x} \in ( \tau_1 (,] \tau_2 ) \wedge
%    \mw{\varphi} ) \leftrightarrow \lnot \forall \msc{x} \lnot ( \msc{x} \in (
%    \tau_1 (,] \tau_2 ) \wedge \mw{\varphi} ) )\label{eq:bl.excd}\tag{bl.excd}
%\end{align}
%\end{restatable}

\begin{proof}\label{proof:bl.excd}
\begin{align}
  1 &&  & \vdash ( \exists \msc{x} ( \msc{x} \in ( \tau_1 (,] \tau_2 ) \wedge
      \mw{\varphi} ) \leftrightarrow \lnot \forall \msc{x}
      \lnot ( \msc{x} \in ( \tau_1 (,] \tau_2 ) \wedge
      \mw{\varphi} ) )&(\mbox{\ref{eq:df-ex}})\notag
\end{align}
\end{proof}

\blexdc
%\begin{restatable}{lemma}{blexdc}~\blTitle{bl.exdc}~ Existential quantification over open-right
%interval [t1 ., t2].
%\begin{align}
%\vdash ( \exists \msc{x} ( \msc{x} \in ( \tau_1 [,) \tau_2 ) \wedge
%    \mw{\varphi} ) \leftrightarrow \lnot \forall \msc{x} \lnot ( \msc{x} \in (
%    \tau_1 [,) \tau_2 ) \wedge \mw{\varphi} ) )\label{eq:bl.exdc}\tag{bl.exdc}
%\end{align}
%\end{restatable}

\begin{proof}\label{proof:bl.exdc}
\begin{align}
  1 &&  & \vdash ( \exists \msc{x} ( \msc{x} \in ( \tau_1 [,) \tau_2 ) \wedge
      \mw{\varphi} ) \leftrightarrow \lnot \forall \msc{x}
      \lnot ( \msc{x} \in ( \tau_1 [,) \tau_2 ) \wedge
      \mw{\varphi} ) )&(\mbox{\ref{eq:df-ex}})\notag
\end{align}
\end{proof}

